\documentclass[12pt, twoside]{article}
%\documentclass[12pt, twoside]{article}
\usepackage[letterpaper, margin=1in, headsep=0.2in]{geometry}
\setlength{\headheight}{0.6in}
%\usepackage[english]{babel}
\usepackage[utf8]{inputenc}
\usepackage{microtype}
\usepackage{amsmath}
\usepackage{amssymb}
%\usepackage{amsfonts}
\usepackage[nomessages]{fp} %\FPeval{\var-name}{2*sin(pi/6)}
\usepackage{siunitx} %units in math. eg 20\milli\meter
\usepackage{yhmath} % for arcs, overparenth command
\usepackage{tikz} %graphics
\usetikzlibrary{quotes, angles, arrows, arrows.meta}
\usepackage{pgfplots}
\usepgfplotslibrary{fillbetween}
\pgfplotsset{compat=1.16}
\usepackage{graphicx} %consider setting \graphicspath{{images/}}
\usepackage{parskip} %no paragraph indent
\usepackage{enumitem}
\usepackage{multicol}
\usepackage{venndiagram}

\usepackage{fancyhdr}
\pagestyle{fancy}
\fancyhf{}
\renewcommand{\headrulewidth}{0pt} % disable the underline of the header
\raggedbottom
\hfuzz=2mm %suppresses overfull box warnings

\usepackage{hyperref}

\fancyhead[LE]{\thepage}
\fancyhead[RO]{Markscheme}
\fancyhead[LO]{La Scuola d'Italia / Dr.\ Huson / IB Math 11 \\* 2 November 2026}

\begin{document}
\subsubsection*{2.7 Classwork: Systems of Equations and Market Equilibrium (no calculator) --- Markscheme}

\begin{enumerate}[itemsep=0.3cm]

%--- Problem 1: Supply and demand lines; estimate equilibrium from graph; solve by elimination; interpret gradient; price ceiling shortage; consumer surplus triangle [13 marks] ---
\item[1.] [Maximum mark: 13]

Demand $D$: $p = 50 - 3q$. Supply $S$: $p = 10 + 2q$. Quantity $q$ in
thousands of cases per month, price $p$ in euros.

\textbf{(a)} $q \approx 8$ (thousand cases) \hfill \textbf{A1}

$p \approx 26$ (euros) \hfill \textbf{A1}

\emph{Accept $7 \le q \le 9$ and $24 \le p \le 28$. Accept the point
written as coordinates, $(8,\ 26)$.}

\textbf{(b)} $3q + p = 50$ and $2q - p = -10$ \hfill \textbf{A1}

\emph{Accept any equivalent pair, for example $-2q + p = 10$. Both
equations must be correct for the A1.}

Adding the equations: $5q = 40$ \hfill \textbf{M1}

$q = 8$ (thousand cases per month) \hfill \textbf{A1}

$p = 50 - 3(8) = 26$ (euros) \hfill \textbf{A1}

\emph{M1 for adding or subtracting the equations (after scaling, if used)
so that one variable is eliminated. Accept eliminating $q$ first:
$2(3q + p) - 3(2q - p)$ gives $5p = 130$. Substitution in place of
elimination earns M0, then A1A1 for correct values. A correct answer with
no working earns A1 for each value only (M0A1A1). FT the final A1 from the
candidate's $q$.}

\textbf{(c)} gradient $= 2$ \hfill \textbf{A1}

For each extra $1000$ cases per month supplied, the price rises by
$2$ euros. \hfill \textbf{R1}

\emph{Accept the reverse reading: for each $2$ euro rise in price, suppliers
offer $1000$ more cases. The R1 needs both the direction (increase) and the
units (euros per thousand cases). ``The line goes up by 2'' with no context
earns R0.}

\textbf{(d)} Substituting $p = 20$ into both equations \hfill \textbf{M1}

Demand: $20 = 50 - 3q$, so $q = 10$ (thousand cases). Supply:
$20 = 10 + 2q$, so $q = 5$ (thousand cases). \hfill \textbf{A1}

Shortage $= 10 - 5 = 5$ thousand cases per month \hfill \textbf{A1}

\emph{The A1 for quantities needs both values. Accept $5000$ cases. A
negative answer, $-5$, or ``a surplus of 5'', earns the final A0. FT the
shortage from the candidate's two quantities, provided demand exceeds
supply.}

\textbf{(e)} Triangle with vertices $(0,\ 50)$, $(0,\ 26)$, $(8,\ 26)$:
base $8$, height $50 - 26 = 24$ \hfill \textbf{M1}

Area $= \frac{1}{2} \times 8 \times 24 = 96$, so the consumer surplus is
$96\,000$ euros \hfill \textbf{A1}

\emph{M1 for the correct base and height (seen in an area calculation),
or for the correct triangle shaded on the graph. Accept $96$ thousand
euros. $96$ euros (units ignored) earns A0. FT from the candidate's
equilibrium point in (b).}

\emph{Wrong methods: the triangle down to the $q$-axis,
$\frac{1}{2} \times 8 \times 26 = 104$, earns M0A0. The rectangle
$8 \times 24 = 192$ earns M0A0. The producer surplus,
$\frac{1}{2} \times 8 \times 16 = 64$, earns M0A0.}

\bigskip \hrule \bigskip

%--- Problem 2: Two-account investment, set up and solve 2x2 by elimination with scaling; extension three-fund 3x3 system [11 marks] ---
\item[2.] [Maximum mark: 11]

\textbf{(a)} $x + y = 10\,000$ \hfill \textbf{A1}

$0.03x + 0.07y = 540$ \hfill \textbf{A1}

\emph{Accept $3x + 7y = 54\,000$ or $\frac{3}{100}x + \frac{7}{100}y = 540$.
$3x + 7y = 540$ (percentages not converted) earns A0.}

\textbf{(b)} $3x + 7y = 54\,000$ and $3x + 3y = 30\,000$ \hfill \textbf{M1}

Subtracting: $4y = 24\,000$, so $y = 6000$ \hfill \textbf{A1}

$x = 10\,000 - 6000 = 4000$ \hfill \textbf{A1}

Giulia puts $4000$ euros in the bond and $6000$ euros in the fund.

\emph{M1 for scaling one or both equations so that a variable can be
eliminated, and adding or subtracting. Substitution earns M0, then
A1A1 for correct values. FT from the candidate's equations in (a). A
correct answer with no working earns A1A1 only.}

\emph{Wrong method: swapping the rates ($0.07x + 0.03y = 540$) gives
$x = 6000$, $y = 4000$; FT from (a) applies only if the swap appears
there, otherwise M1A0A0.}

\textbf{Extension}

\textbf{(c)(i)} \quad
$\begin{aligned}[t]
x + y + z &= 20\,000 \\
0.02x + 0.04y + 0.06z &= 740 \\
y - z &= 2000
\end{aligned}$
\hfill \textbf{A1A1}

\emph{A1 for the interest equation (accept $x + 2y + 3z = 37\,000$ or
$2x + 4y + 6z = 74\,000$). A1 for both the total and the difference
equations ($y = z + 2000$ is fine).}

\textbf{(c)(ii)} Scale the interest equation: $x + 2y + 3z = 37\,000$.

Subtract the total equation: $y + 2z = 17\,000$ \hfill \textbf{M1}

Subtract $y - z = 2000$: $3z = 15\,000$, so $z = 5000$ \hfill \textbf{A1}

$y = 5000 + 2000 = 7000$ \hfill \textbf{A1}

$x = 20\,000 - 7000 - 5000 = 8000$ \hfill \textbf{A1}

\emph{M1 for eliminating one variable to reach a correct equation in the
other two. Any order of elimination is fine. FT from the candidate's
equations in (c)(i). Check: $160 + 280 + 300 = 740$.}

\end{enumerate}
\end{document}
